% ==============================================================
% ARO-KONFORMES REVIEW
% Projektive Strukturtheorie (PST) -- Gesamtreview
% Reviewer: Kimi (Moonshot AI)
% Datum: Juli 2026
% Version: 1.0
% ==============================================================

\documentclass[11pt,a4paper]{article}

\usepackage[utf8]{inputenc}
\usepackage[T1]{fontenc}
\usepackage[german]{babel}
\usepackage{amsmath,amssymb,amsthm}
\usepackage{geometry}
\usepackage{parskip}
\usepackage{enumitem}
\usepackage{booktabs}
\usepackage{xcolor}
\usepackage{hyperref}
\usepackage{longtable}

\geometry{margin=2.5cm}
\setlist{itemsep=0.3em, topsep=0.3em}
\hypersetup{
  colorlinks=true,
  linkcolor=blue,
  urlcolor=blue,
  pdftitle={ARO-Review: PST Gesamtreview},
  pdfauthor={Kimi -- Forschungsprogramm Ontoph\"{a}nomenalismus}
}

\newcommand{\K}[1]{\textbf{[K-#1]}}
\newcommand{\B}[1]{\textbf{[B-#1]}}
\newcommand{\M}[1]{\textbf{[M-#1]}}
\newcommand{\Pos}[1]{\textbf{[P-#1]}}

\newcommand{\Dok}[1]{\textbf{Dokument:} #1\\}
\newcommand{\Stelle}[1]{\textbf{Stelle:} #1\\}
\newcommand{\Prob}[1]{\textbf{Problem:} #1\\}
\newcommand{\Krit}[1]{\textbf{Warum kritisch:} #1\\}
\newcommand{\Empf}[1]{\textbf{Empfehlung:} #1\\}
\newcommand{\Fix}[1]{\textbf{Fixing-Vorschlag:} #1\\}

\title{
  \textbf{Adversarial Review: Projektive Strukturtheorie (PST)}\\[0.4em]
  {\large Gesamtreview PST-1-F bis PST-6-R}\\[0.5em]
  {\normalsize ARO-Konform -- Version~1.0}
}
\author{
  \textbf{Reviewer: Kimi}\\
  \small Forschungsprogramm Ontoph\"{a}nomenalismus
}
\date{Juli 2026}

\begin{document}

\maketitle

\begin{abstract}
Dieses Review bewertet den vollst\"{a}ndigen PST-Dokumentenzyklus
(PST-1-F bis PST-6-R) gem\"{a}\ss{} dem Adversarial Review Protocol (AR).
Es enth\"{a}lt kritische Befunde \textbf{[K]}, bedeutsame Befunde
\textbf{[B]}, minore Befunde \textbf{[M]} und positive Bewertungen
\textbf{[P]}. Alle Befunde sind mit konkreten Fixing-Vorschl\"{a}gen
versehen.
\end{abstract}

\tableofcontents
\newpage

%==============================================================
\section{Zusammenfassung des Gesamteindrucks}
%==============================================================

Die PST ist ein \textbf{methodisch diszipliniertes, philosophisch
koh\"{a}rentes Forschungsprogramm} im Sinne von Lakatos. Der harte Kern
(OPP als pr\"{a}-metrischer Urgrund, Projektion als Erzeugung von
Raumzeit, minimale Ontologie) ist konsistent und anschlussf\"{a}hig an
die philosophischen Module des OP.

Die gr\"{o}\ss{}ten St\"{a}rken liegen in der \textbf{methodischen
Ehrlichkeit} (No-Go-Ergebnisse werden als Ressource genutzt), der
\textbf{Selbstkritik} (Rangneutralit\"{a}t, Zirkularit\"{a}tstests) und
der \textbf{klaren Priorisierung} (PST-5-S).

Die gr\"{o}\ss{}ten Schw\"{a}chen liegen in drei Bereichen:
\begin{enumerate}
  \item \textbf{Mathematische Inkonsistenz:} Der gew\"{a}hlte Weg zur
    Lorentz-Signatur (komplexe Gram-Matrizen, Wick-Rotation,
    imagin\"{a}re Eigenwerte) ist mit der in PST-2-M festgelegten
    reell-symmetrischen Gram-Matrix unvereinbar.
  \item \textbf{Empirische Undokumentiertheit:} Die zentrale
    Behauptung $d_{\text{eff}} \approx 4$ bei $N = 10\,000$ ist
    nicht reproduzierbar, da Code und Rohdaten fehlen.
  \item \textbf{Fehlende philosophische St\"{a}rke:} Die originellste
    Idee des gesamten OP-Programms -- die konstruktiven
    Raumdimensionen vs. die destruktive Zeitdimension als
    ontologische Begr\"{u}ndung der Metrik-Signatur -- taucht in
    keinem PST-Dokument auf.
\end{enumerate}

\newpage

%==============================================================
\section{Kritische Befunde [K]}
%==============================================================

%--------------------------------------------------------------
\subsection{\K{01} PST-1-F: ``Existenz = Bestimmtheit'' ist ein petitio principii}
%--------------------------------------------------------------

\Dok{PST-1-F, Sektion 3}
\Stelle{Theorem ``Existenz = Bestimmtheit''}
\Prob{Der Beweis ist ein Widerspruchsbeweis: Angenommen, X existiert,
aber ist vollkommen unbestimmt. Dann ist X nicht vom Nichts
unterscheidbar. Aber ``existiert'' in der Annahme setzt bereits
Bestimmtheit voraus (per UP).}
\Krit{Dies ist kein Widerspruchsbeweis im strengen Sinne, sondern eine
\textbf{petitio principii}: Die Konklusion ist in die Pr\"{a}misse
eingebaut. Wenn man die UP ablehnt, bricht der Beweis zusammen. Das
Theorem suggeriert einen h\"{o}heren Formalit\"{a}tsgrad als
vorhanden.}
\Empf{Umbenennen in ``Proposition'' oder ``Arbeitshypothese''. Der
Beweis ist g\"{u}ltig \emph{innerhalb} des OP-Systems, aber nicht
absolut.}
\Fix{In PST-1-F, Sektion 3, die Theorem-Umgebung durch eine
\texttt{hypothese}-Umgebung ersetzen und hinzuf\"{u}gen: ``Diese
Proposition ist g\"{u}ltig innerhalb des OP-Axiomensystems (UP als
Axiom). Sie ist kein absoluter Beweis der Gleichsetzung von Existenz
und Bestimmtheit.''}

%--------------------------------------------------------------
\subsection{\K{02} PST-2-M: ``Theorem'' zur Planck-Schranke ist kein Theorem}
%--------------------------------------------------------------

\Dok{PST-2-M, Sektion 7}
\Stelle{Theorem ``Planck-Schranke f\"{u}r bestimmte Existenz''}
\Prob{Das Dokument pr\"{a}sentiert eine Behauptung als
\texttt{theorem}-Umgebung:
\[
  \mathcal{B}(X,Y) = 1 \iff MI(X,Y) \geq I_0 \iff d_{XY} \geq \ell_P
\]
Die Methodennote sagt im selben Atemzug: ``Dieses Theorem ist eine
Arbeitshypothese ... noch nicht mathematisch bewiesen.''}
\Krit{In LaTeX (und in der Mathematik) ist eine \texttt{theorem}-Umgebung
eine \textbf{Behauptung mit Beweis}. Eine Arbeitshypothese geh\"{o}rt in
eine \texttt{hypothese}- oder \texttt{conjecture}-Umgebung. Die Mischung
von \texttt{theorem} mit ``ist eine Arbeitshypothese'' ist ein
\textbf{Selbstwiderspruch auf der Pr\"{a}sentationsebene}. Ein externer
Mathematiker w\"{u}rde hier sofort aufh\"{o}ren zu lesen.}
\Empf{In zwei Schritte aufteilen: (1) \texttt{lemma} oder
\texttt{proposition}: $MI(X,Y) \geq I_0 \iff \mathcal{B}(X,Y) = 1$
(Definition/Operationalisierung). (2) \texttt{hypothese}: Diese
Bestimmtheit korrespondiert mit $d_{XY} \geq \ell_P$
(physikalische Br\"{u}cke).}
\Fix{In PST-2-M, Sektion 7, die \texttt{theorem}-Umgebung streichen und
durch zwei separate Umgebungen ersetzen: Eine
\texttt{definition}/\texttt{proposition} f\"{u}r die
Informationsbedingung und eine \texttt{hypothese} f\"{u}r die
Planck-Korrespondenz.}

%--------------------------------------------------------------
\subsection{\K{03} PST-2-M: Die Verbindung zwischen \textbf{R} und \textbf{K} fehlt}
%--------------------------------------------------------------

\Dok{PST-2-M, Sektion 3 und 4}
\Stelle{Definition von \textbf{R} und Definition von \textbf{K}}
\Prob{Es werden zwei Matrizen definiert:
\begin{itemize}
  \item $\mathbf{R} \in \mathbb{R}^{N \times N}$ mit $R_{ij} = r(s_i, s_j)$,
    wobei $r: S \times S \to \mathbb{R}_{\geq 0}$ eine ``Relationsst\"{a}rke'' ist.
  \item $\mathbf{K} \in \mathbb{R}^{N \times N}$ mit $K_{ij} = k(s_i, s_j)$,
    wobei $k$ ein ``symmetrisch-positiv semidefiniter Kernel'' ist.
\end{itemize}
Aber: \textbf{Wie h\"{a}ngen $\mathbf{R}$ und $\mathbf{K}$ zusammen?}}
\Krit{Die Hypothese in Sektion 3 suggeriert: $r(s_i, s_j) = MI(X_i; X_j)$.
Aber in Sektion 4 wird $K_{ij} = k(s_i, s_j)$ definiert, ohne Bezug zu
$\mathbf{R}$. Wenn $\mathbf{K} \neq \mathbf{R}$, dann ist die gesamte
Kette ``Relation $\to$ Kernel $\to$ Spektrum $\to$ Geometrie''
unterbrochen. Wenn $\mathbf{K} = \mathbf{R}$, dann muss $\mathbf{R}$
symmetrisch-positiv semidefinit sein -- was bei einer allgemeinen
Relationsst\"{a}rke $r: S \times S \to \mathbb{R}_{\geq 0}$ \textbf{nicht
gegeben} ist.}
\Empf{Explizit definieren: Entweder $\mathbf{K} = \mathbf{R}$ und
$\mathbf{R}$ muss als Kernel konstruiert werden, oder $\mathbf{R}$ ist
die Rohdaten-Matrix und $\mathbf{K} = \varphi(\mathbf{R})$ ist eine
Kernel-Konstruktion.}
\Fix{In PST-2-M, Sektion 4, vor der Definition von $\mathbf{K}$ eine
Zwischendefinition einf\"{u}gen: ``Die Gram-Matrix $\mathbf{K}$ wird aus
der Relationsmatrix $\mathbf{R}$ durch einen positiv-semidefiniten
Kernel konstruiert: $K_{ij} = \exp(-R_{ij}^2 / 2\sigma^2)$
(Gau\ss{}scher Kernel) oder $K_{ij} = R_{ij}$ falls $\mathbf{R}$ bereits
positiv semidefinit ist.''}

%--------------------------------------------------------------
\subsection{\K{04} PST-2-M: ``Isometrische Einbettung'' ist mathematisch irref\"{u}hrend}
%--------------------------------------------------------------

\Dok{PST-2-M, Sektion 4}
\Stelle{Theorem ``Spektralzerlegung der Gram-Matrix''}
\Prob{Das Theorem behauptet:
\[
  \Phi: s_i \mapsto (\sqrt{\lambda_1} v_{1i}, \ldots, \sqrt{\lambda_N} v_{Ni})
\]
sei ``eine isometrische Einbettung der Zust\"{a}nde in den euklidischen
Raum $\mathbb{R}^N$''.}
\Krit{Eine isometrische Einbettung erh\"{a}lt Abst\"{a}nde. Aber:
\begin{itemize}
  \item Der OPP hat \textbf{keine Metrik} (per Definition).
  \item Die Einbettung $\Phi$ ist die Kernel-PCA-Einbettung. Sie erh\"{a}lt
    das Skalarprodukt im Feature-Raum (definiert durch $\mathbf{K}$),
    aber nicht einen ``Abstand'' im OPP -- denn einen solchen gibt es
    nicht.
\end{itemize}
Das Wort ``isometrisch'' suggeriert, dass die Einbettung eine
vorgegebene Metrik erh\"{a}lt. Aber genau das lehnt die PST ab. Es ist
ein \textbf{Kategorienfehler}: Die PST behauptet, Metrik entstehe erst
durch die Einbettung -- aber dann kann die Einbettung nicht
``isometrisch'' sein.}
\Empf{Formulieren als: ``Die Einbettung $\Phi$ induziert ein
Skalarprodukt und damit eine Metrik auf dem Bild der Zust\"{a}nde in
$\mathbb{R}^N$. Diese Metrik ist emergent -- sie existiert nicht im
OPP, sondern erst in der projizierten Darstellung.''}
\Fix{In PST-2-M, Sektion 4, den Satz ``eine isometrische Einbettung
...'' ersetzen durch: ``Die Einbettung $\Phi$ repr\"{a}sentiert die
durch $\mathbf{K}$ definierte \"{A}hnlichkeitsstruktur im euklidischen
Raum $\mathbb{R}^N$ und induziert dort eine emergente Metrik.''}

%--------------------------------------------------------------
\subsection{\K{05} PST-3-P: Die Stabilit\"{a}tshypothese f\"{u}r $d = 4$ ist zirkul\"{a}r}
%--------------------------------------------------------------

\Dok{PST-3-P, Sektion 4}
\Stelle{Hypothese ``Stabilit\"{a}tshypothese f\"{u}r $d = 4$''}
\Prob{Die Hypothese begr\"{u}ndet $d = 4$ mit physikalischen Argumenten
aus der \emph{bekannten} Physik:
\begin{itemize}
  \item Planetenbahnen sind nur in $d = 3$ stabil
  \item Yang-Mills-Theorien sind in $d = 4$ renormierbar
  \item $SO(1,3)$ hat in $d = 4$ Weyl-Spinoren
\end{itemize}}
\Krit{Die PST soll die Physik \emph{aus dem OPP herleiten}, nicht die
bekannte Physik als Begr\"{u}ndung f\"{u}r $d = 4$ verwenden. Wenn ihr
sagt: ``$d = 4$ emergiert, weil in $d = 4$ Planetenbahnen stabil
sind'', dann setzt ihr voraus, was ihr beweisen wollt. Das ist ein
\textbf{petitio principii} auf der Ebene der Physik.}
\Empf{Entweder (a) die physikalischen Argumente als \emph{heuristische
Randbedingungen} kennzeichnen (``Wir wissen aus der bekannten Physik,
dass $d = 4$ funktioniert; die PST muss zeigen, \emph{warum} der OPP
notwendigerweise $d = 4$ liefert'') -- oder (b) die Argumente durch
\emph{strukturelle} Argumente ersetzen, die unabh\"{a}ngig von der
bekannten Physik sind (z.B. aus der Spektraltheorie der Gram-Matrix,
aus der Stabilit\"{a}t von Diffusion Maps, aus der
Informationsgeometrie).}
\Fix{In PST-3-P, Sektion 4, die Hypothese umbenennen in
``Heuristische Randbedingung'' und formulieren: ``Die bekannte Physik
liefert die Randbedingung $d = 4$ f\"{u}r physikalische
Projektionsoperatoren. Die PST sucht einen Operator, der diese
Randbedingung erf\"{u}llt, ohne sie vorauszusetzen.''}

%--------------------------------------------------------------
\subsection{\K{06} PST-3-P: Die ``Definition'' von $\mathcal{P}^*$ ist eine Wunschliste}
%--------------------------------------------------------------

\Dok{PST-3-P, Sektion 5}
\Stelle{Definition ``Universeller Projektionsoperator''}
\Prob{Ein mathematisches Objekt wird definiert durch seine
\emph{Eigenschaften}, nicht durch seine \emph{Ziele}. Die
``Definition'' listet f\"{u}nf Bedingungen auf, die $\mathcal{P}^*$
erf\"{u}llen soll -- aber sie sagt nicht, \emph{was} $\mathcal{P}^*$
ist. Ist es eine Abbildung? Ein Funktor? Ein Operator auf einem
Hilbert-Raum? Ein stochastischer Kern?}
\Krit{Ohne mathematische Spezifikation ist $\mathcal{P}^*$ ein
Platzhalter, kein Objekt. Das ist f\"{u}r ein Forschungsprogramm in
Ordnung -- aber es darf nicht als \texttt{definition} in einem
mathematischen Dokument erscheinen.}
\Empf{Umbenennen in \texttt{hypothese} oder \texttt{desideratum}. Die
f\"{u}nf Bedingungen sind \emph{Randbedingungen}, keine Definition.}
\Fix{In PST-3-P, Sektion 5, die \texttt{definition}-Umgebung durch
\texttt{hypothese} oder ein neues \texttt{desideratum} ersetzen. Die
f\"{u}nf Bedingungen als ``Randbedingungen f\"{u}r $\mathcal{P}^*$''
kennzeichnen. Hinzufigen: ``Eine vollst\"{a}ndige mathematische
Definition von $\mathcal{P}^*$ erfordert die in PST-5-S, Priorit\"{a}t 3
geforderte Fundierung (Kategorientheorie, Operatoralgebren oder
spektrale Tripel).''}

%--------------------------------------------------------------
\subsection{\K{07} PST-3-P: Lorentz-Signatur aus ``Imagin\"{a}rteil von Eigenvektoren'' ist mathematisch problematisch}
%--------------------------------------------------------------

\Dok{PST-3-P, Sektion 7}
\Stelle{Hypothese ``Lorentz-Signatur aus Projektion''}
\Prob{Die Hypothese suggeriert, dass die Minkowski-Signatur
$(+,-,-,-)$ aus dem Imagin\"{a}rteil von Eigenvektoren bei komplexen
Kernel-Matrizen emergieren k\"{o}nnte. Aber:
\begin{itemize}
  \item In PST-2-M ist die Gram-Matrix \textbf{reell-symmetrisch}
    definiert: $K_{ij} = k(s_i, s_j) \in \mathbb{R}$.
  \item Reell-symmetrische Matrizen haben \textbf{reelle Eigenwerte und
    reelle Eigenvektoren}. Es gibt keinen Imagin\"{a}rteil.
  \item Wenn man zu komplexen Kerneln \"{u}bergeht, verliert man die
    Symmetrie $\mathbf{K} = \mathbf{K}^T$ und damit die Garantie
    reeller Eigenwerte. Dann ist die Spektralzerlegung nicht mehr
    garantiert reell.
\end{itemize}}
\Krit{Das ist ein \textbf{mathematischer Kategorienfehler}. Man kann
nicht aus einer reellen Struktur einen Imagin\"{a}rteil zaubern, der
dann die Lorentz-Signatur erzeugt. Das w\"{a}re, als wollte man aus einem
reellen Vektorraum eine komplexe Struktur extrahieren -- das geht
nur, wenn man eine komplexe Struktur \emph{voraussetzt}.}
\Empf{Entweder (a) die Gram-Matrix von vorneherein als hermitesch
(komplex) definieren -- was die gesamte PST-2-M \"{a}ndern w\"{u}rde --
oder (b) einen anderen Mechanismus f\"{u}r die Lorentz-Signatur suchen.
Ein vielversprechender Ansatz w\"{a}re: Die Signatur emergiert aus der
\textbf{Ordnung} der Eigenwerte und deren \emph{Vorzeichenstruktur}
in einer indefiniten Metrik (z.B. durch eine Pseudometrik mit
Signatur $(1,3)$, die aus der Gram-Matrix durch eine nicht-triviale
Bilinearform induziert wird). Aber das erfordert einen indefiniten
Kernel -- und der ist in PST-2-M nicht definiert.}
\Fix{In PST-3-P, Sektion 7, die Hypothese \"{u}berarbeiten oder
streichen. Als Ersatz eine neue Hypothese einf\"{u}hren, die auf der
\textbf{konstruktiv/destruktiv-Idee} basiert (siehe Fixing-Vorschlag
zu \textbf{[K-15]}). Die alte Hypothese als ``verworfener Ansatz''
in eine Methodennote verlagern.}

%--------------------------------------------------------------
\subsection{\K{08} PST-4-G: Die zentrale Behauptung $d_{\text{eff}} \approx 4$ ist nicht verifizierbar}
%--------------------------------------------------------------

\Dok{PST-4-G, Sektion 5}
\Stelle{Befund ``Zentraler PST-Befund: $d_{\text{eff}} \approx 4$''}
\Prob{Das Dokument behauptet:
\begin{quote}
``Diffusion Maps mit $t \approx 9$--$10$ liefern reproduzierbar:
$d_{\text{eff}} = 3.98 \pm 0.36$ \"{u}ber verschiedene OPP-Stichproben
($N = 1\,000$ bis $N = 100\,000$), verschiedene Kernel-Funktionen,
verschiedene $K$-Werte ($K = 10$ bis $K = 50$), verschiedene
Schwellenwerte $\theta$ ($0.90$ bis $0.99$).''
\end{quote}
Aber: \textbf{Es gibt keinen Code, keine Rohdaten, keine
GitHub-Links.} Die einzige verifizierbare Simulation, die wir bisher
gesehen haben, ist \textbf{Test 35} ($N = 40$, 2D-Punkte, 10 Runs),
bei dem $d_{\text{eff}}$ bei $t = 10$ exakt bei $3.98$ lag -- aber als
Teil eines \textbf{monotonen Abfalls} ($t = 6 \to 7.8$, $t = 10 \to
4.0$, $t = 12 \to 3.1$).}
\Krit{Ein wissenschaftliches Dokument, das reproduzierbare Ergebnisse
beansprucht, muss die \textbf{Reproduktionsgrundlage} liefern. Ohne
Code und Rohdaten ist das ein \textbf{Glaubensbekenntnis}, kein
Befund. Ein externer Reviewer (oder eine zuk\"{u}nftige KI) kann diese
Behauptung nicht pr\"{u}fen.}
\Empf{Entweder (a) den Code und die Rohdaten sofort auf GitHub
ver\"{o}ffentlichen (mit Commit-Hash, Zeitstempel, verwendeter
Hardware) -- oder (b) die Behauptung auf ``Vorl\"{a}ufiger Befund aus
internen Tests, Dokumentation folgt'' abschw\"{a}chen.}
\Fix{In PST-4-G, Sektion 5, einen neuen Absatz einf\"{u}gen:
\begin{verbatim}
\begin{methodennote}[Reproduzierbarkeit und Datenlage]
Die hier berichteten Ergebnisse stammen aus internen
Testreihen, die derzeit nicht vollst\"{a}ndig versioniert sind.
Eine vollst\"{a}ndige Dokumentation inklusive Code, Rohdaten,
Parameter-Metadaten und Reproduktions-Skripten ist in
Vorbereitung (siehe Framework-Dokumentation, geplant f\"{u}r
PST-7-F). Bis zur Ver\"{o}ffentlichung dieser Daten gilt der
Befund als vorl\"{a}ufig.
\end{methodennote}
\end{verbatim}}

%--------------------------------------------------------------
\subsection{\K{09} PST-4-G: Die Monotonie von $d_{\text{eff}}(t)$ wird nicht thematisiert}
%--------------------------------------------------------------

\Dok{PST-4-G, Sektion 5}
\Stelle{Remark ``Bedeutung von $t \approx 9$--$10$''}
\Prob{Das Dokument beschreibt:
\begin{itemize}
  \item Kleines $t$: $d_{\text{eff}}$ streut stark
  \item Mittleres $t$ (9--10): $d_{\text{eff}}$ konvergiert stabil auf $\approx 4$
  \item Gro\ss{}es $t$ (20--50): $d_{\text{eff}}$ nimmt wieder ab
\end{itemize}
Aber es \textbf{zeigt nicht die vollst\"{a}ndige Kurve}. Wir wissen aus
Test 35, dass $d_{\text{eff}}$ monoton mit $t$ sinkt. Die
``Konvergenz bei $t \approx 9$--$10$'' k\"{o}nnte ein \textbf{Plateau}
in diesem Abfall sein -- oder ein \textbf{Artefakt} der spezifischen
Parameter ($N$, Kernel, $\theta$).}
\Krit{Wenn $d_{\text{eff}}(t)$ eine glatte, monotone Funktion ist,
dann ist der Wert 4 bei $t = 10$ nicht fundamental. Er ist ein
\textbf{Schnitt durch eine kontinuierliche Kurve}. Die Physik erfordert
aber eine \textbf{fundamentale} Dimension -- nicht eine, die von einem
freien Parameter $t$ abh\"{a}ngt.}
\Empf{Die vollst\"{a}ndige Kurve $d_{\text{eff}}(t)$ f\"{u}r
verschiedene $N$ und Kernel dokumentieren. Wenn es ein Plateau gibt,
zeigen. Wenn es keines gibt, ehrlich sagen: ``Der Wert 4 ist
parameterabh\"{a}ngig und nicht fundamental.''}
\Fix{In PST-4-G, Sektion 5, die Tabelle der Test-35-Ergebnisse
vollst\"{a}ndig einf\"{u}gen (alle $t$-Werte von 6 bis 12 mit
Standardabweichungen). Hinzufigen: ``Die Monotonie von
$d_{\text{eff}}(t)$ zeigt, dass der Wert 4 kein fundamentales
Merkmal des OPP ist, sondern eine Eigenschaft des
Diffusion-Map-Operators bei bestimmten Parametern. Die Frage,
ob es ein physikalisch ausgezeichnetes $t$ gibt, bleibt offen.''}

%--------------------------------------------------------------
\subsection{\K{10} PST-4-G: $N = 100\,000$ ist unglaubw\"{u}rdig ohne Kontext}
%--------------------------------------------------------------

\Dok{PST-4-G, Sektion 5.2}
\Stelle{Befund ``Reproduzierbarkeit''}
\Prob{Die Behauptung $N = 100\,000$ wirft Fragen auf:
\begin{itemize}
  \item $N = 100\,000$ Zust\"{a}nde mit einer dichten Gram-Matrix
    erfordern \textbf{80 GB RAM} ($100\,000^2 \times 8$ Byte).
  \item Die Eigenwertberechnung einer $100\,000 \times 100\,000$
    Matrix ist \textbf{rechenintensiv} ($O(N^3)$, also $\sim 10^{15}$
    Operationen).
  \item Auf einem i7-10700K w\"{u}rde das \textbf{Tage} dauern.
  \item Der Autor hat gesagt, die Rechner sind noch nicht aufgebaut.
\end{itemize}}
\Krit{Die Behauptung $N = 100\,000$ ohne Kontext ist eine
\textbf{Glaubensbehauptung}. Wurde das auf Grok ausgef\"{u}hrt? (Grok
hat Zeitlimits von Minuten, nicht Tagen.) Wurde das auf Cloud-Hardware
ausgef\"{u}hrt? (Dann: Welche? Kosten?) Oder ist $N = 100\,000$ eine
\textbf{Projektion} f\"{u}r zuk\"{u}nftige Tests?}
\Empf{Explizit benennen, welche $N$ auf welcher Hardware mit welcher
Laufzeit getestet wurden. ``$N = 100\,000$'' ohne Kontext ist eine
Behauptung, die das Vertrauen untergr\"{a}bt.}
\Fix{In PST-4-G, Sektion 5.2, eine Tabelle einf\"{u}gen:
\begin{verbatim}
\begin{center}
\begin{tabular}{llll}
\toprule
\textbf{N} & \textbf{Hardware} & \textbf{Laufzeit} & \textbf{Status} \\
\midrule
40         & Grok (Cloud)      & $<$ 1 Min.        & Dokumentiert \\
1\,000     & [noch offen]      & [noch offen]      & Geplant \\
10\,000    & [noch offen]      & [noch offen]      & Geplant \\
100\,000   & i7-10700K + 256GB & gesch\"{a}tzt 3--5 Tage & Zuk\"{u}nftig \\
\bottomrule
\end{tabular}
\end{center}
\end{verbatim}}

%--------------------------------------------------------------
\subsection{\K{11} PST-5-S: Der ``harte Kern'' ist zu hart}
%--------------------------------------------------------------

\Dok{PST-5-S, Sektion 1}
\Stelle{Definition ``Harter Kern der PST''}
\Prob{Der vierte Punkt des harten Kerns lautet: ``Es gibt einen
universellen Projektionsoperator $\mathcal{P}^*$.'' Das ist eine
\textbf{massive Annahme}. Wenn sich herausstellt, dass es keinen
\emph{einen} universellen Operator gibt, sondern eine \textbf{Klasse}
von kontextabh\"{a}ngigen Operatoren (je nach Beobachter, Energieskala,
etc.), w\"{u}rde das die PST nicht zwingend scheitern lassen. Es
w\"{u}rde nur bedeuten, dass $\mathcal{P}^*$ kein einzelner Operator,
sondern ein \textbf{Funktor} oder eine \textbf{Familie} ist.}
\Krit{In Lakatos' Terminologie geh\"{o}rt in den harten Kern nur das,
was bei Aufgabe das Programm beendet. Die Existenz eines \emph{einzigen}
universellen Operators ist st\"{a}rker als n\"{o}tig. Die Aussage
``Raumzeit emergiert durch Projektion'' ist hart. Die Aussage ``Es
gibt genau ein $\mathcal{P}^*$'' ist eine Hilfshypothese im
Schutzg\"{u}rtel.}
\Empf{Punkt 3 des harten Kerns \"{a}ndern zu: ``Es gibt eine Klasse von
Projektionsoperatoren, die aus dem OPP Raumzeit emergieren lassen.''
Und $\mathcal{P}^*$ als spezifische Hypothese in den Schutzg\"{u}rtel
verschieben.}
\Fix{In PST-5-S, Sektion 1, die Definition des harten Kerns
\"{a}ndern:
\begin{verbatim}
\begin{definition}[Harter Kern der PST -- korrigiert]
\begin{enumerate}
  \item Der OPP existiert als hochdimensionaler, pr\"{a}-metrischer
    Zustandsraum.
  \item Raumzeit, Physik und Bestimmtheit emergieren durch Projektion.
  \item Es gibt eine Klasse von Projektionsoperatoren, die aus dem OPP
    Raumzeit emergieren lassen.
  \item Das Prinzip der minimalen Ontologie gilt.
\end{enumerate}
\end{definition}
\end{verbatim}}

%--------------------------------------------------------------
\subsection{\K{12} PST-5-S und PST-6-R: Der Weg zur Lorentz-Signatur \"{u}ber
komplexe Gram-Matrizen ist eine Sackgasse}
%--------------------------------------------------------------

\Dok{PST-5-S, Sektion 4 (Priorit\"{a}t 2); PST-6-R, Sektion 2}
\Stelle{``Imagin\"{a}re Eigenwerte bei komplexen Gram-Matrizen'' und
``Wick-Rotation''}
\Prob{Dieser Ansatz taucht in drei Dokumenten auf (PST-3-P, PST-5-S,
PST-6-R) und ist mathematisch mit der in PST-2-M festgelegten
reell-symmetrischen Gram-Matrix \textbf{unvereinbar}.
\begin{itemize}
  \item Eine komplexwertige Gram-Matrix verliert die Garantie
    reeller Eigenwerte und positiver Semidefinitheit.
  \item Die Wick-Rotation in der QFT setzt \textbf{bereits} die
    Minkowski-Metrik voraus -- sie ist ein Trick \emph{innerhalb}
    der Physik, nicht eine Herleitung \emph{der} Physik.
  \item Der Ansatz verfestigt sich in der Forschungsagenda, obwohl
    er ein mathematischer Kategorienfehler ist.
\end{itemize}}
\Krit{Die PST verfolgt einen mathematisch unhaltbaren Ansatz und
ignoriert gleichzeitig eine vielversprechendere Alternative (siehe
\textbf{[B-01]}). Das ist eine \textbf{Fehlallokation von
Aufmerksamkeit} mit potenziell fatalen Konsequenzen f\"{u}r die
weitere Entwicklung.}
\Empf{Diesen Ansatz aus der Priorit\"{a}tenliste streichen oder als
``hochspekulativ, erfordert Neufundierung der Gram-Matrix-Theorie''
kennzeichnen. Stattdessen den Weg \"{u}ber die
\textbf{konstruktiv/destruktive Dimension} (Raum = konstruktiv, Zeit
= destruktiv) als prim\"{a}ren Kandidaten nennen.}
\Fix{In PST-5-S, Sektion 4, Priorit\"{a}t 2 ersetzen durch:
\begin{verbatim}
\begin{ziel}[Priorit\"{a}t 2: Lorentz-Signatur -- korrigiert]
\textbf{Frage:} Wie emergiert $(+,-,-,-)$ aus einer signaturfreien
Gram-Matrix?

\textbf{Kandidatenansatz (neu):}
Die drei Raumdimensionen sind ``konstruktiv'' -- sie st\"{a}rken die
Bildung von Attraktoren. Die Zeitdimension ist ``destruktiv''
-- sie zerst\"{o}rt Attraktoren (Entropie, Informationsverlust).
Die Signatur $(+,-,-,-)$ reflektiert diesen ontologischen Gegensatz.
Mathematisch k\"{o}nnte dies durch eine indefinite Bilinearform
auf der Gram-Matrix realisiert werden, bei der konstruktive Eigenwerte
positiv und destruktive Eigenwerte negativ sind.

\textbf{Verworfener Ansatz:} Imagin\"{a}re Eigenwerte bei komplexen
Gram-Matrizen (inkonsistent mit PST-2-M; siehe Review [K-12]).
\end{ziel}
\end{verbatim}}

%--------------------------------------------------------------
\subsection{\K{13} PST-6-R: Wick-Rotation ist physikalisch irref\"{u}hrend}
%--------------------------------------------------------------

\Dok{PST-6-R, Sektion 2}
\Stelle{Hypothese ``Wick-Rotation als projektive Eigenschaft''}
\Prob{Die Wick-Rotation $t \to -i\tau$ in der QFT ist ein
\textbf{mathematisches Hilfsmittel}, keine physikalische
Transformation. Sie macht das Pfadintegral konvergent, indem sie
von Minkowski- zu Euklidischer Zeit wechselt. Aber: Die QFT setzt
\textbf{bereits} die Minkowski-Metrik voraus.}
\Krit{Die PST behauptet, die Lorentz-Signatur aus dem OPP
herzuleiten. Wenn sie daf\"{u}r die Wick-Rotation verwendet, setzt sie
implizit voraus, was sie beweisen will: die Existenz einer
Zeitkoordinate $t$, die man rotieren kann. Das ist ein
\textbf{petitio principii}.}
\Empf{Streichen oder als ``heuristische Analogie, keine Herleitung''
kennzeichnen. Die Wick-Rotation kann nicht das Fundament sein, auf dem
die Lorentz-Signatur emergiert.}
\Fix{In PST-6-R, Sektion 2, die Hypothese umbenennen in
\texttt{remark} mit dem Text: ``Die Wick-Rotation in der QFT ist ein
Standardtrick zur Konvergenz des Pfadintegrals. Sie setzt jedoch
bereits die Minkowski-Metrik voraus und kann daher nicht als
Herleitung der Lorentz-Signatur aus dem OPP dienen. Sie bleibt eine
heuristische Analogie.''}

%--------------------------------------------------------------
\subsection{\K{14} PST-6-R: Die Schr\"{o}dinger-Gleichung als
``projektive Eigenschaft'' ist zu vage}
%--------------------------------------------------------------

\Dok{PST-6-R, Sektion 4}
\Stelle{Hypothese ``Schr\"{o}dinger-Gleichung aus Projektion''}
\Prob{Die Hypothese sagt: ``Die Schr\"{o}dinger-Gleichung k\"{o}nnte
eine projektive Eigenschaft sein.'' Aber sie liefert keinen
Mechanismus. Die Schr\"{o}dinger-Gleichung enth\"{a}lt:
\begin{itemize}
  \item Eine Ableitung nach der Zeit ($\partial/\partial t$)
  \item Einen Hamilton-Operator $\hat{H}$
  \item Die imagin\"{a}re Einheit $i$
\end{itemize}
Keines dieser Elemente ist in der PST bisher definiert.}
\Krit{Die Hypothese suggeriert Substanz, wo keine ist. Sie ist nicht
falsch -- aber sie ist leer. Das ist f\"{u}r ein Forschungsprogramm in
Ordnung, aber die Sprache sollte transparenter sein.}
\Empf{Entweder (a) explizit als ``noch v\"{o}llig offen, keine
mathematische Substanz'' kennzeichnen, oder (b) einen konkreten
Ansatz skizzieren.}
\Fix{In PST-6-R, Sektion 4, die Hypothese durch eine
\texttt{methodennote} ersetzen:
\begin{verbatim}
\begin{methodennote}[Schr\"{o}dinger-Gleichung -- offenes Terrain]
Die Herleitung der Schr\"{o}dinger-Gleichung aus der Projektion
ist ein Fernziel der PST. Bisher existiert kein mathematischer
Mechanismus, der die Elemente der Gleichung ($i$, $\hat{H}$,
$\partial/\partial t$) aus PST-Strukturen ableitet. Dieser
Abschnitt kartografiert das Terrain, nicht den Weg.
\end{methodennote}
\end{verbatim}}

\newpage

%==============================================================
\section{Bedeutsame Befunde [B]}
%==============================================================

%--------------------------------------------------------------
\subsection{\B{01} PST-1-F bis PST-6-R: Die destruktive Zeitdimension
fehlt komplett}
%--------------------------------------------------------------

\Dok{Alle PST-Dokumente}
\Stelle{Fehlend}
\Prob{Die st\"{a}rkste originelle Idee des gesamten OP-Programms --
die konstruktiven Raumdimensionen vs. die destruktive Zeitdimension
als ontologische Begr\"{u}ndung der Metrik-Signatur -- taucht in
keinem PST-Dokument auf.}
\Krit{Die PST verfolgt einen mathematischen Weg (komplexe
Eigenwerte, Wick-Rotation), der problematisch ist, w\"{a}hrend sie
einen philosophisch fundierten Weg (konstruktiv/destruktiv)
ignoriert, der mathematisch sauberer w\"{a}re. Das ist ein
\textbf{Blindfleck}.}
\Empf{Diese Idee muss in eine \"{U}berarbeitung der PST
einflie\ss{}en. Sie k\"{o}nnte die Lorentz-Signatur ontologisch
begr\"{u}nden, ohne die mathematische Konsistenz der Gram-Matrix zu
zerst\"{o}ren.}
\Fix{In PST-5-S, Sektion 4, als neue Priorit\"{a}t einf\"{u}gen:
\begin{verbatim}
\begin{ziel}[Priorit\"{a}t 2b: Ontologische Begr\"{u}ndung der
Lorentz-Signatur]
\textbf{Hypothese:} Die drei Raumdimensionen sind ``konstruktiv''
-- sie st\"{a}rken die Bildung und Stabilit\"{a}t von Attraktoren.
Die Zeitdimension ist ``destruktiv'' -- sie zerst\"{o}rt Attraktoren
(Entropie, Informationsverlust). Die Signatur $(+,-,-,-)$
reflektiert diesen ontologischen Gegensatz.

\textbf{Mathematische Umsetzung:} Eine indefinite Bilinearform
auf der Gram-Matrix, bei der konstruktive Eigenwerte
positiv und destruktive Eigenwerte negativ sind. Dies w\"{u}rde
eine indefinite Metrik mit Signatur $(1,3)$ oder $(3,1)$ erzeugen,
ohne die reelle Struktur der Gram-Matrix aufzugeben.
\end{ziel}
\end{verbatim}
In PST-6-R, Sektion 2, eine neue Hypothese einf\"{u}hren, die diesen
Ansatz entwickelt.}

%--------------------------------------------------------------
\subsection{\B{02} PST-3-P: Die Verbindung $I_0 \leftrightarrow \hbar$
bleibt dimensionsanalytisch offen}
%--------------------------------------------------------------

\Dok{PST-3-P, Sektion 6}
\Stelle{Hypothese ``Quantisierung aus endlicher Projektionsaufl\"{o}sung''}
\Prob{$[I_0] = \text{Bit}$ (Informationsma\ss{}), $[\hbar] =
\text{J} \cdot \text{s} = \text{kg} \cdot \text{m}^2/\text{s}$
(Wirkung). Die Hypothese behauptet eine Korrespondenz, aber ohne
eine \textbf{Konstante}, die Bits in Wirkung umrechnet, ist das eine
Typengleichung, keine physikalische Gleichung.}
\Krit{In der Physik ist jede Gleichung dimensionshomogen. Wenn
$I_0 = \hbar$ gelten soll, braucht man einen Faktor $k$ mit
$[k] = \text{J} \cdot \text{s}/\text{Bit}$. Woher kommt dieser
Faktor? Aus dem OPP? Aus der Projektion? Das ist die entscheidende
Frage.}
\Empf{Explizit als offene Frage formulieren: ``Die Korrespondenz
$I_0 \leftrightarrow \hbar$ erfordert eine dimensionsanalytische
Br\"{u}cke, die derzeit nicht vorhanden ist.''}
\Fix{In PST-3-P, Sektion 6, die Hypothese erg\"{a}nzen:
\begin{verbatim}
\begin{methodennote}[Dimensionsanalytische L\"{u}cke]
Die Korrespondenz $I_0 \leftrightarrow \hbar$ erfordert eine
Br\"{u}ckenkonstante $k$ mit $[k] = \text{J} \cdot
\text{s}/\text{Bit}$. Ein m\"{o}glicher Ansatz: $I_0 = \hbar /
(k_B T \cdot \tau)$, wobei $T$ und $\tau$ emergente Gr\"{o}\ss{}en
der Projektion w\"{a}ren. Dies bleibt ein offenes Problem.
\end{methodennote}
\end{verbatim}}

%--------------------------------------------------------------
\subsection{\B{03} PST-3-P: Die ``Proposition'' Beobachter =
Projektionsoperatoren ist ein Postulat}
%--------------------------------------------------------------

\Dok{PST-3-P, Sektion 8}
\Stelle{Proposition ``Beobachter = Projektionsoperatoren''}
\Prob{Eine Proposition in der Mathematik ist eine Behauptung, die
aus Definitionen und Axiomen \emph{bewiesen} werden kann. Hier wird
behauptet: ``Ein physikalischer Beobachter ist identisch mit einem
Projektionsoperator.'' Das ist keine Folgerung -- das ist eine
\textbf{Setzung}, ein Postulat.}
\Krit{Die Proposition suggeriert Beweisbarkeit, wo keine ist. Das
ist ein Versto\ss{} gegen die MKO-Sprachregel (Ebene 1 statt Ebene
3/5).}
\Empf{Umbenennen in \texttt{postulat} oder \texttt{definition}. Oder:
Als \texttt{proposition} belassen, aber mit einem Beweis versehen,
der aus den Axiomen des OP folgt.}
\Fix{In PST-3-P, Sektion 8, die \texttt{proposition} durch
\texttt{postulat} ersetzen und hinzuf\"{u}gen: ``Dies ist eine
grundlegende Setzung der PST, keine ableitbare Proposition. Sie
verbindet die philosophische Kategorie des Beobachters (DPR) mit
der mathematischen Kategorie des Projektionsoperators.''}

%--------------------------------------------------------------
\subsection{\B{04} PST-4-G: Die Traum-Analogie / SSD-Analogie fehlt}
%--------------------------------------------------------------

\Dok{PST-4-G, Gesamtdokument}
\Stelle{Fehlend}
\Prob{Das Dokument beschreibt den Paradigmenwechsel von ``OPP hat
Dimensionen'' zu ``Projektion erzeugt Dimensionen''. Aber es
vermisst die \textbf{intuitive Verankerung}, die der Autor dem
Reviewer erkl\"{a}rt hat: Die SSD-Analogie (Bits $\to$ Player $\to$
Film).}
\Krit{Die SSD-Analogie ist die beste didaktische Br\"{u}cke zwischen
Philosophie und Mathematik. Ihr Fehlen macht PST-4-G f\"{u}r externe
Leser schwerer zug\"{a}nglich.}
\Empf{Einen Abschnitt ``Intuitive Verankerung'' hinzuf\"{u}gen.}
\Fix{In PST-4-G, nach Sektion 2 (Phase 2), eine neue Sektion
einf\"{u}gen:
\begin{verbatim}
\begin{remark}[Intuitive Verankerung: Die SSD-Analogie]
Der OPP ist wie eine SSD: statisch, zeitlos, nur Bits.
Der Diffusion-Map-Operator ist wie ein Player: Er projiziert
sequentiell und erzeugt emergent Raum und Zeit.
$d_{\text{eff}} \approx 4$ ist die emergente Dimension des Films.
Diese Analogie ist keine mathematische Herleitung, sondern ein
didaktisches Hilfsmittel (MKO Ebene 4).
\end{remark}
\end{verbatim}}

%--------------------------------------------------------------
\subsection{\B{05} PST-5-S: Die Verbindung zu Connes ist kontr\"{a}r,
nicht komplement\"{a}r}
%--------------------------------------------------------------

\Dok{PST-5-S, Sektion 7}
\Stelle{``PST und nicht-kommutative Geometrie (Connes)''}
\Prob{Connes' spektrale Tripel $(A, H, D)$ setzt einen
\textbf{Dirac-Operator $D$} voraus, der eine \textbf{Metrik} kodiert.
Die PST behauptet, Metrik emergiere erst durch Projektion. Das sind
\textbf{kontr\"{a}re Ausgangspunkte}, nicht komplement\"{a}re.}
\Krit{Wenn die PST sagt ``Metrik ist emergent'', dann kann sie
nicht gleichzeitig Connes' Ansatz als ``eine mathematisch
ausgearbeitete Version des PST-Projektionsprinzips'' interpretieren.
Connes \emph{setzt} die Metrik durch $D$ voraus; die PST
\emph{leitet} sie ab.}
\Empf{Die Verbindung zu Connes als ``interessante mathematische
Analogie, aber mit umgekehrter ontologischer Richtung''
kennzeichnen -- oder streichen.}
\Fix{In PST-5-S, Sektion 7, den Text \"{a}ndern:
\begin{verbatim}
\begin{remark}[PST und nicht-kommutative Geometrie (Connes)]
Alain Connes beschreibt Raumzeit \"{u}ber spektrale Tripel
$(A, H, D)$. Der Dirac-Operator $D$ kodiert dabei bereits eine
Metrik. Die PST und Connes' Ansatz sind daher ontologisch
kontr\"{a}r: Connes setzt Metrik voraus, die PST leitet sie ab.
Ob beide Programme letztlich formal \"{a}quivalent sind, ist eine
offene Frage. Die PST kann Connes' Mathematik als Werkzeug
nutzen, nicht als ontologische Grundlage.
\end{remark}
\end{verbatim}}

%--------------------------------------------------------------
\subsection{\B{06} PST-5-S: Das Scheitern-Kriterium 3 ist mathematisch trivial}
%--------------------------------------------------------------

\Dok{PST-5-S, Sektion 5}
\Stelle{``Wann scheitert die PST?'' -- Kriterium 3}
\Prob{``Die Lorentz-Signatur grunds\"{a}tzlich nicht aus einer reellen,
positiv semidefiniten Gram-Matrix emergieren kann.'' Das ist kein
m\"{o}gliches Scheitern -- das ist ein \textbf{mathematisches
Faktum}. Eine positiv semidefinite Matrix hat nur nicht-negative
Eigenwerte. Eine indefinite Metrik mit Signatur $(+,-,-,-)$ hat
positive und negative Eigenwerte. Das eine kann nicht aus dem
anderen folgen.}
\Krit{Das Scheitern-Kriterium suggeriert, dass dies eine offene
empirische Frage ist. Aber es ist eine geschlossene mathematische.
Die PST muss entweder (a) indefinite Gram-Matrizen zulassen, (b)
die Signatur durch einen anderen Mechanismus erzeugen, oder (c)
zugeben, dass die reell-positive Definitheit eine zu starke Annahme
war.}
\Empf{Das Kriterium umformulieren: ``Es zeigt sich, dass keine
Modifikation der Gram-Matrix-Struktur (z.B. \"{U}bergang zu indefiniten
Bilinearformen) die Lorentz-Signatur konsistent erzeugen kann.''}
\Fix{In PST-5-S, Sektion 5, Kriterium 3 ersetzen durch:
\begin{verbatim}
\item Es zeigt sich, dass keine konsistente Modifikation der
  Gram-Matrix-Struktur (einschlie\ss{}lich indefiniter
  Bilinearformen) die Lorentz-Signatur $(+,-,-,-)$ erzeugen
  kann.
\end{verbatim}}

%--------------------------------------------------------------
\subsection{\B{07} PST-6-R: Die Verbindung zur ART ist zu oberfl\"{a}chlich}
%--------------------------------------------------------------

\Dok{PST-6-R, Sektion 6}
\Stelle{``Allgemeine Relativit\"{a}tstheorie (ART)''}
\Prob{Die Behauptung, die Einstein-Feldgleichungen ``k\"{o}nnten aus
Eigenschaften des Projektionsspektrums emergieren'', ist nicht
falsch -- aber sie ist \textbf{leer}. Es wird nicht erkl\"{a}rt:
\begin{itemize}
  \item Wie $G_{\mu\nu}$ (Einstein-Tensor) aus einer Gram-Matrix folgen
    soll
  \item Wie $T_{\mu\nu}$ (Energie-Impuls-Tensor) als
    Projektionseigenschaft interpretiert werden soll
  \item Wie die Kovarianz (allgemeine Koordinatentransformationen)
    emergieren soll
\end{itemize}}
\Krit{Die Sprache suggeriert mehr Substanz als vorhanden. Das ist
im Rahmen eines Forschungsprogramms legitim -- aber die Sprache
sollte transparenter sein zwischen ``intuitive Analogie'' und
``mathematische Hypothese''.}
\Empf{Entweder einen konkreten mathematischen Ansatz skizzieren
oder als ``v\"{o}llig offen'' kennzeichnen.}
\Fix{In PST-6-R, Sektion 6, den Text \"{a}ndern:
\begin{verbatim}
\begin{remark}[Allgemeine Relativit\"{a}tstheorie (ART) -- offen]
Die Verbindung zwischen PST und ART ist v\"{o}llig offen.
Insbesondere fehlt jeder mathematische Mechanismus, der:
\begin{itemize}
  \item den Einstein-Tensor $G_{\mu\nu}$ aus einer Gram-Matrix
    ableitet,
  \item den Energie-Impuls-Tensor $T_{\mu\nu}$ als
    Projektionseigenschaft interpretiert,
  \item die allgemeine Kovarianz aus der Projektionsstruktur
    emergieren l\"{a}sst.
\end{itemize}
Dieser Abschnitt kartografiert das Terrain, nicht den Weg.
\end{remark}
\end{verbatim}}

\newpage

%==============================================================
\section{Minore Befunde [M]}
%==============================================================

%--------------------------------------------------------------
\subsection{\M{01} PST-2-M: Die ``Koh\"{a}renzerhaltung'' ist nicht
mathematisch definiert}
%--------------------------------------------------------------

\Dok{PST-2-M, Sektion 4; PST-3-P, Sektion 2}
\Stelle{Definition ``Physikalische Projektion'', Punkt 1 und 3}
\Prob{``Globale Koh\"{a}renzerhaltung'' und ``Koh\"{a}renzerhaltung''
sind nicht formal definiert. Was bedeutet ``Koh\"{a}renz"
mathematisch? Ist es eine Bedingung an den Kernel? An die
Eigenwerte? An die Diffusion?}
\Fix{In PST-2-M oder PST-3-P eine formale Definition einf\"{u}gen:
``Eine Projektion $\mathcal{P}$ erh\"{a}lt die Koh\"{a}renz, wenn f\"{u}r
alle Zust\"{a}nde $s_i, s_j$ gilt: $R_{ij} > \varepsilon \Rightarrow
d_{\mathcal{O}}(\mathcal{P}(s_i), \mathcal{P}(s_j)) < \delta$'' --
oder \"{a}hnlich. Ohne formale Definition bleibt ``Koh\"{a}renz'' ein
intuitiver Begriff.}

%--------------------------------------------------------------
\subsection{\M{02} PST-2-M: OPP als $(S, \mathcal{R})$ ist zu allgemein}
%--------------------------------------------------------------

\Dok{PST-2-M, Sektion 2}
\Stelle{Definition ``OPP als metrisch-freier Zustandsraum''}
\Prob{$\mathcal{M}_{\text{OPP}} = (S, \mathcal{R})$ mit $S$ als
``Menge von Zust\"{a}nden'' und $\mathcal{R}$ als ``Kollektion von
Relationen'' ist mathematisch fast leer. Eine ``Kollektion von
Relationen'' ist keine strukturierte mathematische Kategorie.}
\Fix{Den OPP entweder als \textbf{Kategorie} definieren (Objekte =
Zust\"{a}nde, Morphismen = Relationen mit Kompositionsregeln) oder als
\textbf{Graph} (Knoten = Zust\"{a}nde, Kanten = Relationsst\"{a}rken).
Beides w\"{a}re mathematisch pr\"{a}ziser als ``Kollektion''.}

%--------------------------------------------------------------
\subsection{\M{03} PST-2-M: ``Die Felddynamik bildet $\mathcal{A}$ auf
sich selbst ab'' -- Kontamination?}
%--------------------------------------------------------------

\Dok{PST-2-M, Sektion 2}
\Stelle{Definition ``Attraktoren als stabile Zust\"{a}nde'', Punkt
``Invarianz''}
\Prob{``Die Felddynamik bildet $\mathcal{A}$ auf sich selbst ab.''
Das ist derselbe Satz, der bereits in OAD (v1.3) und MGO (v1.0)
stand und dort die Kritik am Zeit-Paradoxon ausgel\"{o}st hat. Hier,
in PST-2-M, steht er in der Definition des OPP -- dem
mathematisch pr\"{a}-metrischen Raum.}
\Fix{Ersetzen durch: ``Es existiert eine strukturerhaltende
Abbildung $f: \mathcal{A} \to \mathcal{A}$ mit $f(\mathcal{A}) =
\mathcal{A}$.'' Oder, kategorientheoretisch: ``$\mathcal{A}$ ist
invariant unter den Endomorphismen des OPP.''}

%--------------------------------------------------------------
\subsection{\M{04} PST-3-P: Die Verbindung zur Kaluza-Klein-Theorie
ist zu vage}
%--------------------------------------------------------------

\Dok{PST-3-P, Sektion 3}
\Stelle{Remark ``Verbindung zu Kaluza-Klein''}
\Prob{Die Bemerkung sagt: ``Es gibt eine strukturelle \"{A}hnlichkeit
zur Kaluza-Klein-Theorie.'' Aber es wird nicht erkl\"{a}rt, \emph{was}
genau \"{a}hnlich ist. In KK werden Eichfelder als
Metrik-Komponenten in zus\"{a}tzlichen Dimensionen interpretiert. In
der PST werden Eichgruppen als \emph{verbleibende Symmetrien} einer
Projektion interpretiert. Das sind zwei v\"{o}llig verschiedene
Mechanismen.}
\Fix{Entweder pr\"{a}zisieren: ``Beide Theorien leiten Eichstrukturen
aus einer h\"{o}herdimensionalen Geometrie ab; KK durch
Kompaktifizierung, PST durch Projektion'' -- oder streichen, wenn
die \"{A}hnlichkeit zu oberfl\"{a}chlich ist.}

\newpage

%==============================================================
\section{Positive Bewertungen [P]}
%==============================================================

\Pos{01} \textbf{PST-1-F: Die Kontaminationsregel als mathematisches
Arbeitsprinzip.} Die Forderung, keine Metrik, Topologie oder
Dimension als Voraussetzung zu verwenden, ist die konsequente
Umsetzung der philosophischen Pr\"{a}misse und ein Alleinstellungsmerkmal.

\Pos{02} \textbf{PST-2-M: Die Kernel-Methods-Strategie ist mathematisch
solide.} Die Verwendung von Gram-Matrizen, Spektralzerlegung und
Diffusion Maps ist etabliert und m\"{a}chtig. Die Idee, $d_{\text{eff}}$
aus dem Spektrum abzulesen, ist der Standardweg in der
Dimensionsreduktion.

\Pos{03} \textbf{PST-2-M: Die Ehrlichkeit bez\"{u}glich des Status.}
Die Methodennoten (MKO Ebene 4--5) sind konsequent. Das Dokument
behauptet nicht mehr, als es liefern kann.

\Pos{04} \textbf{PST-3-P: Die Drei-Ebenen-Struktur (Ontisch /
Projektiv / Ph\"{a}nomenal).} Das ist die konsequente Umsetzung dessen,
was DPR philosophisch entwickelt hat. Die Tabelle in Sektion 1 ist
das R\"{u}ckgrat der gesamten PST.

\Pos{05} \textbf{PST-3-P: Projektion als Symmetriebrechung.} Statt zu
fragen ``Warum hat der OPP Symmetrien?'' fragt ihr: ``Welche
Symmetrien \emph{bleiben \"{u}brig}, wenn der OPP projiziert wird?''
Das ist genau der richtige Weg f\"{u}r eine Emergenztheorie.

\Pos{06} \textbf{PST-4-G: Die Dokumentation negativer Ergebnisse als
wissenschaftliche Ressource.} Das No-Go-Ergebnis ($d_{\text{eff}}^{\text{roh}}
\approx 30$--$45$) wird nicht verdr\"{a}ngt, sondern als Fundament
f\"{u}r den Paradigmenwechsel genutzt. Der Vergleich mit Michelson-Morley
ist nicht \"{u}berzogen.

\Pos{07} \textbf{PST-4-G: Die Selbstkritik zum $K$-Parameter und zum
Stabilit\"{a}tsfunktional.} Die Erkenntnis, dass $\lambda_F$ implizit
niedrige Dimensionen bevorzugt, und die Konsequenz
(``Rangneutralit\"{a}t als Pflicht'') sind professionell.

\Pos{08} \textbf{PST-5-S: Der Lakatos-Rahmen.} Sich selbst als
``Forschungsprogramm im Sinne von Lakatos'' zu verstehen -- mit
hartem Kern, Schutzg\"{u}rtel und Heuristik -- entzieht Kritikern die
M\"{o}glichkeit, einzelne Fehler als Todessto\ss{} zu verkaufen.

\Pos{09} \textbf{PST-5-S: Die Falsifizierbarkeit.} Die
Scheitern-Kriterien in Sektion 5 sind Popper-konform und pr\"{u}fbar.

\Pos{10} \textbf{PST-6-R: Die Dynamik-als-Beobachterpfad-Hypothese.}
Physikalische Dynamik ist kein Prozess im OPP selbst, sondern die
Sequenz von Projektionen, die ein Beobachter-Attraktor durchl\"{a}uft.
Das ist der philosophisch st\"{a}rkste Gedanke des gesamten PST-Zyklus.

\Pos{11} \textbf{PST-6-R: Die Kausalit\"{a}t-als-Projektionsreihenfolge-Hypothese.}
Ursache vor Wirkung = Projektion von $t_1$ reicht in $t_2$ hinein,
nicht umgekehrt. Das gibt der Zeitrichtung eine informationale
Begr\"{u}ndung ohne thermodynamische Voraussetzungen.

\Pos{12} \textbf{PST-6-R: Die Ehrlichkeit bez\"{u}glich der
Naturkonstanten.} ``Die PST macht bisher keine Vorhersage \"{u}ber
Naturkonstanten.'' Das ist der einzige Weg, wie man mit dem
Feinabstimmungsproblem umgehen kann -- ohne zu l\"{u}gen.

\newpage

%==============================================================
\section{\"{U}bergeordnete systematische Punkte}
%==============================================================

\subsection{Die drei gr\"{o}\ss{}ten systematischen Probleme der PST}

\begin{enumerate}
  \item \textbf{Der mathematische Weg zur Lorentz-Signatur ist eine
    Sackgasse.} Die wiederholte Verwendung von ``komplexen
    Gram-Matrizen'' und ``Wick-Rotation'' in PST-3-P, PST-5-S und
    PST-6-R zeigt, dass das Team an einem mathematisch unhaltbaren
    Ansatz festh\"{a}lt. Die Gram-Matrix in PST-2-M ist
    reell-symmetrisch und positiv semidefinit. Diese Struktur kann
    keine indefinite Metrik mit Signatur $(+,-,-,-)$ erzeugen. Der
    Weg muss entweder durch indefinite Gram-Matrizen oder durch die
    konstruktiv/destruktiv-Idee ersetzt werden.

  \item \textbf{Die empirische Basis ist nicht dokumentiert.} Die
    zentrale Behauptung ($d_{\text{eff}} \approx 4$ bei $N = 10\,000$)
    ist nicht reproduzierbar. Code und Daten fehlen. Das untergr\"{a}bt
    die Glaubw\"{u}rdigkeit des gesamten numerischen Arms der PST.

  \item \textbf{Die Sprache suggeriert mehr Substanz als vorhanden.}
    Viele Hypothesen verwenden physikalische Begriffe
    (Schr\"{o}dinger-Gleichung, Einstein-Feldgleichungen,
    Eichgruppen) ohne mathematische Substanz. Das ist im Rahmen eines
    Forschungsprogramms legitim -- aber die Sprache sollte
    transparenter sein zwischen ``intuitiver Analogie'' und
    ``mathematischer Hypothese''.
\end{enumerate}

\subsection{Die drei gr\"{o}\ss{}ten St\"{a}rken der PST}

\begin{enumerate}
  \item \textbf{Die Pr\"{a}zisierung der Frage.} Von ``Hat der OPP 4D?''
    zu ``Welche Projektionen erzeugen 4D?'' -- das ist der
    entscheidende methodische Schritt.

  \item \textbf{Die methodische Disziplin.} Rangneutralit\"{a}t,
    Selbstkritik, Falsifizierbarkeit -- das ist auf einem Niveau,
    das selbst etablierte theoretische Physik selten erreicht.

  \item \textbf{Die philosophische Koh\"{a}renz.} Die PST bleibt
    konsequent im OP-Rahmen. Sie vermeidet die Trennung von
    Philosophie und Mathematik, die viele andere Programme plagt.
\end{enumerate}

\newpage

%==============================================================
\section{Priorisierte Empfehlungen f\"{u}r die \"{U}berarbeitung}
%==============================================================

\begin{center}
\begin{tabular}{clp{8cm}l}
\toprule
\textbf{Priorit\"{a}t} & \textbf{Dokument} & \textbf{Problem} & \textbf{Aufwand} \\
\midrule
1 & PST-4-G & Code und Daten ver\"{o}ffentlichen; $d_{\text{eff}} \approx 4$
  verifizierbar machen & Hoch \\
2 & PST-5-S & Harten Kern korrigieren ($\mathcal{P}^*$ als Klasse, nicht
  Einzeloperator) & Niedrig \\
3 & PST-3-P, PST-5-S, PST-6-R & Lorentz-Signatur-Weg \"{u}berarbeiten:
  komplexe Eigenwerte/Wick-Rotation streichen, konstruktiv/destruktiv
  einf\"{u}hren & Hoch \\
4 & PST-2-M & Verbindung $\mathbf{R} \to \mathbf{K}$ definieren & Mittel \\
5 & PST-2-M & ``Isometrische Einbettung'' korrigieren & Niedrig \\
6 & PST-1-F & ``Existenz = Bestimmtheit'' als Hypothese kennzeichnen & Niedrig \\
7 & PST-3-P & $\mathcal{P}^*$ als Hypothese/Desideratum kennzeichnen & Niedrig \\
8 & PST-4-G & SSD/Traum-Analogie als intuitive Verankerung einf\"{u}gen & Niedrig \\
9 & PST-3-P & $I_0 \leftrightarrow \hbar$: Dimensionsl\"{u}cke benennen & Mittel \\
10 & PST-6-R & ART und Schr\"{o}dinger als ``offenes Terrain'' kennzeichnen & Niedrig \\
\bottomrule
\end{tabular}
\end{center}

\vspace{1em}

\noindent\textbf{Wichtigste Empfehlung:} Bevor neue Simulationen
(PST-7-F) oder neue mathematische Ans\"{a}tze entwickelt werden,
m\"{u}ssen die drei kritischen Punkte [K-08], [K-09], [K-12]
adressiert werden. Ohne reproduzierbare Daten und ohne einen
konsistenten Weg zur Lorentz-Signatur riskiert die PST, in
Spekulation abzugleiten.

\vspace{2em}

\begin{center}
\rule{0.5\textwidth}{0.4pt}
\end{center}

\vspace{1em}

\noindent\textit{Dieses Review wurde gem\"{a}\ss{} dem Adversarial Review
Protocol (AR) verfasst. Der Reviewer hat die Probleme benannt und
L\"{o}sungsvorschl\"{a}ge skizziert, aber die inhaltliche Entscheidung
und Umsetzung obliegt dem Autorenteam.}

\end{document}

